\section*{Chapter \ref{ch:g11:dissolution} --- Dissolving Ionic and Molecular Solids}

\begin{solution}{exo:g11:dissolution:1}
\ce{NaCl(s) -> Na+(aq) + Cl-(aq)};
\ce{CaCl2(s) -> Ca^{2+}(aq) + 2Cl-(aq)};
\ce{Na2SO4(s) -> 2Na+(aq) + SO4^{2-}(aq)};
\ce{AlCl3(s) -> Al^{3+}(aq) + 3Cl-(aq)}.
\end{solution}

\begin{solution}{exo:g11:dissolution:2}
$[\ce{Na+}] = 2 \times 0.050 = \qty{0.10}{mol/L}$;
$[\ce{SO4^{2-}}] = \qty{0.050}{mol/L}$.
\end{solution}

\begin{solution}{exo:g11:dissolution:3}
\omterm{def:g11:dissolution:solvation}{Hydration} is the surrounding of the \omterm{def:g9:ions:ion}{ion} by water \omterm{def:g7:atoms-and-molecules:molecule}{molecules} held to it.
The oxygen end ($\delta^-$) points towards a \omterm{def:g9:ions:ion}{cation}, the hydrogen end
($\delta^+$) towards an \omterm{def:g9:ions:ion}{anion}: opposite charges attract.
\end{solution}

\begin{solution}{exo:g11:dissolution:4}
$M(\ce{AlCl3}) = 27.0 + 3 \times 35.5 = \qty{133.5}{g/mol}$;
$n = 2.67 / 133.5 = \qty{0.0200}{mol}$; $c = 0.0200 / 0.2000 =
\qty{0.100}{mol/L}$; $[\ce{Al^{3+}}] = \qty{0.100}{mol/L}$,
$[\ce{Cl-}] = \qty{0.300}{mol/L}$.
\end{solution}

\begin{solution}{exo:g11:dissolution:5}
The head \ce{-COO-} is ionic, attracted by water: \omterm{def:g11:dissolution:hydrophilic}{hydrophilic}. The tail
of 17 carbon \omterm{def:g7:atoms-and-molecules:atom}{atoms}, with only \ce{C-C} and nearly non-polar \ce{C-H}
bonds, is non-polar: \omterm{def:g11:dissolution:hydrophilic}{hydrophobic}.
\end{solution}

\begin{solution}{exo:g11:dissolution:6}
Salt: water (ionic solid, polar \omterm{def:g6:solutions-and-solubility:solute}{solvent}). Wax: cyclohexane (non-polar
chains). Sugar: water (its \ce{O-H} groups form \omterm{def:g11:polarity-and-cohesion:hydrogen-bond}{hydrogen bonds} with
water). Iodine: cyclohexane (\omterm{def:g11:polarity-and-cohesion:polar-molecule}{non-polar molecule}).
\end{solution}

\begin{solution}{exo:g11:dissolution:7}
The \ce{O-H} group of ethanol forms \omterm{def:g11:polarity-and-cohesion:hydrogen-bond}{hydrogen bonds} with water \omterm{def:g7:atoms-and-molecules:molecule}{molecules},
which replace those it breaks: ethanol fits into water. Hexane, non-polar,
can form no such bonds; the water \omterm{def:g7:atoms-and-molecules:molecule}{molecules}, held to one another by
\omterm{def:g11:polarity-and-cohesion:hydrogen-bond}{hydrogen bonds}, do not let it in.
\end{solution}

\begin{solution}{exo:g11:dissolution:8}
The \omterm{def:g11:dissolution:hydrophilic}{hydrophobic} tails avoid the water and \omterm{def:g2:mixing-and-dissolving:dissolve}{dissolve} in the grease, at the
centre; the charged heads are attracted by water and stay outside. All
the \omterm{def:g11:dissolution:amphiphilic}{micelles} carry negative charges on their surfaces, and like charges
repel: they stay apart.
\end{solution}

\begin{solution}{exo:g11:dissolution:9}
Cyclohexane: it is \omterm{def:g6:solutions-and-solubility:miscible}{immiscible} with water, so it forms a separate layer
that can be run off, and the scent (non-polar) \omterm{def:g2:mixing-and-dissolving:dissolve}{dissolves} well in it.
Ethanol is \omterm{def:g6:solutions-and-solubility:miscible}{miscible} with water: no second layer forms, and nothing can be
separated.
\end{solution}

\begin{solution}{exo:g11:dissolution:10}
Cyclohexane is less dense than water. The lower, aqueous layer is run
off through the tap. The iodine is in the upper cyclohexane layer, which
stays in the funnel.
\end{solution}

\begin{solution}{exo:g11:dissolution:11}
Nothing: copper sulfate is ionic and does not \omterm{def:g2:mixing-and-dissolving:dissolve}{dissolve} in non-polar
cyclohexane. The water stays blue and the cyclohexane colourless.
\end{solution}

\begin{solution}{exo:g11:dissolution:12}
$[\ce{Cl-}] = 2c$, so $c = \qty{0.150}{mol/L}$;
$n = 0.150 \times 0.2500 = \qty{0.0375}{mol}$;
$m = 0.0375 \times 111.1 = \qty{4.17}{g}$.
\end{solution}

\begin{solution}{exo:g11:dissolution:13}
The calcium and magnesium \omterm{def:g9:ions:ion}{ions} \omterm{def:g9:ions:precipitate}{precipitate} the stearate \omterm{def:g9:ions:ion}{ions} as scum:
the soap is used up before it can form \omterm{def:g11:dissolution:amphiphilic}{micelles}, and lathers badly.
Calcium alone: \qty{0.010}{mol} of \ce{Ca^{2+}} takes
\qty{0.020}{mol} of stearate, that is $0.020 \times 306.0 =
\qty{6.1}{g}$ of sodium stearate per litre.
\end{solution}

\begin{solution}{exo:g11:dissolution:14}
Volume \qty{0.2000}{L}. \ce{Na+}: $0.020 / 0.2000 = \qty{0.10}{mol/L}$;
\ce{Ca^{2+}}: $0.010 / 0.2000 = \qty{0.050}{mol/L}$; \ce{Cl-}:
$(0.020 + 0.020) / 0.2000 = \qty{0.20}{mol/L}$. Positive charge:
$0.10 + 2 \times 0.050 = 0.20$; negative: $0.20$. Neutral.
\end{solution}

\begin{solution}{exo:g11:dissolution:15}
A litre of heptane has a mass of \qty{0.68}{kg} and can hold
$17 \times 0.68 = \qty{12}{g}$ of iodine, about 40 times more than a
litre of water. Shaken together, nearly all the iodine passes into the
heptane: the \omterm{def:g10:chemical-species:extraction}{extraction} is efficient.
\end{solution}

\begin{solution}{pb:g11:dissolution:1}
\textbf{1.} From the rocks the water has run through (limestone,
dolomite, gypsum), which \omterm{def:g2:mixing-and-dissolving:dissolve}{dissolve} a little.

\textbf{2.} $[\ce{Ca^{2+}}] = 0.120 / 40.1 = \qty{2.99e-3}{mol/L}$;
$[\ce{Mg^{2+}}] = 0.024 / 24.3 = \qty{9.9e-4}{mol/L}$.

\textbf{3.} \ce{Ca(HCO3)2 -> Ca^{2+}(aq) + 2HCO3-(aq)}; so
$[\ce{HCO3-}] = 2 \times \num{2.99e-3} = \qty{5.99e-3}{mol/L}$.

\textbf{4.} $\num{2.99e-3} \times 150 = \qty{0.449}{mol}$.

\textbf{5.} $18 \times 12.0 + 35 \times 1.0 + 2 \times 16.0 + 23.0 =
\qty{306.0}{g/mol}$.

\textbf{6.} \ce{C17H35COONa(s) -> C17H35COO-(aq) + Na+(aq)}.

\textbf{7.} Head \ce{-COO-}: \omterm{def:g11:dissolution:hydrophilic}{hydrophilic}; tail \ce{C17H35-}:
\omterm{def:g11:dissolution:hydrophilic}{hydrophobic}. It has both kinds of part: \omterm{def:g11:dissolution:amphiphilic}{amphiphilic}.

\textbf{8.} A tiny sphere of soap \omterm{def:g9:ions:ion}{ions}, tails inwards, heads outwards.
The tails \omterm{def:g2:mixing-and-dissolving:dissolve}{dissolve} in the grease, which is broken into droplets wrapped
in \omterm{def:g11:dissolution:amphiphilic}{micelles}; the charged surfaces keep the droplets apart in the water,
which rinses them away.

\textbf{9.} The soap must form \omterm{def:g9:ions:ion}{ions} and \omterm{def:g11:dissolution:amphiphilic}{micelles}, which needs water
around the heads; and the grease is carried off by the rinsing water. In
oil the grease would simply stay dissolved on the skin.

\textbf{10.} \ce{2C17H35COO- + Ca^{2+} -> Ca(C17H35COO)2}: charges
$2 \times (-1) + 2 = 0$ on the left, $0$ on the right.

\textbf{11.} \ce{C17H35COO-}: excess $- 2x$; \ce{Ca^{2+}}:
$\num{2.99e-3} - x$; calcium stearate: $x$. The calcium \omterm{def:g9:ions:ion}{ions} are
limiting: $x_{\max} = \qty{2.99e-3}{mol}$.

\textbf{12.} $2 \times \num{2.99e-3} = \qty{5.99e-3}{mol}$ per litre.

\textbf{13.} $M = 36 \times 12.0 + 70 \times 1.0 + 4 \times 16.0 + 40.1 =
\qty{606.1}{g/mol}$; $\num{2.99e-3} \times 606.1 = \qty{1.81}{g}$ of scum
per litre.

\textbf{14.} $2 \times \num{9.9e-4} = \qty{1.98e-3}{mol}$ per litre.

\textbf{15.} $2 \times 0.449 = \qty{0.898}{mol}$.

\textbf{16.} $0.898 \times 306.0 = \qty{275}{g}$.

\textbf{17.} Magnesium: $\num{9.9e-4} \times 150 = \qty{0.148}{mol}$,
which takes \qty{0.296}{mol} of stearate, \qty{91}{g} of soap; in all
about $275 + 91 = \qty{366}{g}$.

\textbf{18.} $275 / 100 \approx 2.7$ bars.

\textbf{19.} The calcium of one bath wastes about \qty{0.27}{kg} of soap
as scum.
\end{solution}
